\documentclass{article}
\usepackage{amsmath, amssymb, amsthm}

\title{International Competition in Mathematics for University Students\\
Plovdiv, Bulgaria, 1995}
\date{}

\begin{document}
\maketitle

\begin{center}\textbf{First day}\end{center}

\section*{Problem 1 (Day 1)}
Let $X$ be a nonsingular matrix with columns $X_1, X_2, \ldots, X_n$. Let $Y$ be a matrix with columns $X_2, X_3, \ldots, X_n, 0$. Show that $A = YX^{-1}$ and $B = X^{-1}Y$ have rank $n-1$ and only zero eigenvalues.

\subsection*{Solution}
Let $J = (a_{ij})$ be the $n\times n$ matrix with $a_{ij} = 1$ if $i = j+1$ and $0$ otherwise (subdiagonal shift). Then $\mathrm{rank}(J) = n-1$ and its only eigenvalues are $0$. Moreover $Y = XJ$, so
\[
A = YX^{-1} = XJX^{-1}, \qquad B = X^{-1}Y = J.
\]
Both are similar to $J$ (or equal to $J$), so both have rank $n-1$ with only zero eigenvalues.

\section*{Problem 2 (Day 1)}
Let $f$ be continuous on $[0,1]$ with $\int_x^1 f(t)\,dt \geq \frac{1-x^2}{2}$ for every $x \in [0,1]$. Show $\int_0^1 f^2(t)\,dt \geq \frac{1}{3}$.

\subsection*{Solution}
From $0 \leq \int_0^1 (f(x) - x)^2 dx = \int_0^1 f^2 - 2\int_0^1 xf + \int_0^1 x^2$ we get
\[
\int_0^1 f^2(x)\,dx \geq 2\int_0^1 xf(x)\,dx - \tfrac{1}{3}.
\]
Interchanging order of integration,
\[
\int_0^1\!\int_x^1 f(t)\,dt\,dx = \int_0^1 t f(t)\,dt \geq \int_0^1 \frac{1-x^2}{2}\,dx = \frac{1}{3}.
\]
Thus $\int_0^1 f^2 \geq 2 \cdot \tfrac{1}{3} - \tfrac{1}{3} = \tfrac{1}{3}$.

\section*{Problem 3 (Day 1)}
Let $f$ be twice continuously differentiable on $(0, +\infty)$ with $\lim_{x\to 0+} f'(x) = -\infty$ and $\lim_{x\to 0+} f''(x) = +\infty$. Show
\[
\lim_{x\to 0+} \frac{f(x)}{f'(x)} = 0.
\]

\subsection*{Solution}
There is $r > 0$ with $f'(x) < 0$ and $f''(x) > 0$ on $(0, r)$. So $f$ is decreasing and $f'$ is increasing on $(0,r)$. By MVT, for $0 < x < x_0 < r$:
\[
f(x) - f(x_0) = f'(\xi)(x - x_0) > 0, \quad \xi \in (x, x_0).
\]
Since $f'(x) < f'(\xi) < 0$,
\[
x - x_0 < \frac{f'(\xi)}{f'(x)}(x - x_0) = \frac{f(x) - f(x_0)}{f'(x)} < 0.
\]
Taking $x \to 0+$:
\[
-x_0 \leq \liminf_{x\to 0+} \frac{f(x)}{f'(x)} \leq \limsup_{x\to 0+} \frac{f(x)}{f'(x)} \leq 0.
\]
This holds for every $x_0 \in (0, r)$, hence the limit equals $0$.

\section*{Problem 4 (Day 1)}
Let $F: (1, \infty) \to \mathbb{R}$ be $F(x) = \int_x^{x^2} \frac{dt}{\ln t}$. Show $F$ is injective and find its range.

\subsection*{Solution}
We compute
\[
F'(x) = \frac{2x}{\ln x^2} - \frac{1}{\ln x} = \frac{x}{\ln x} - \frac{1}{\ln x} = \frac{x-1}{\ln x} > 0 \quad (x > 1).
\]
So $F$ is strictly increasing, hence one-to-one. Since
\[
F(x) \geq (x^2 - x) \min_{x \leq t \leq x^2} \frac{1}{\ln t} = \frac{x^2 - x}{\ln x^2} \to \infty
\]
as $x \to \infty$, the range is $(F(1+), \infty)$. Substituting $t = e^v$:
\[
F(x) = \int_{\ln x}^{2\ln x} \frac{e^v}{v}\,dv,
\]
so $x \ln 2 < F(x) < x^2 \ln 2$. Hence $F(1+) = \ln 2$ and the range is $(\ln 2, \infty)$.

\section*{Problem 5 (Day 1)}
Let $A, B$ be real $n \times n$ matrices. Suppose there exist $n+1$ distinct reals $t_1, \ldots, t_{n+1}$ such that each $C_i = A + t_i B$ is nilpotent. Show $A$ and $B$ are nilpotent.

\subsection*{Solution}
We expand $(A + tB)^n = A^n + tP_1 + t^2 P_2 + \cdots + t^{n-1} P_{n-1} + t^n B^n$ where $P_i$ depend only on $A,B$. Fix $(i,j)$ indices; let $a, p_1, \ldots, p_{n-1}, b$ be those entries. The polynomial
\[
bt^n + p_{n-1}t^{n-1} + \cdots + p_1 t + a
\]
has at least $n+1$ roots $t_1, \ldots, t_{n+1}$. All coefficients vanish, so $A^n = 0$, $B^n = 0$, and $P_i = 0$; hence $A, B$ are nilpotent.

\section*{Problem 6 (Day 1)}
Let $p > 1$. Show there is $K_p > 0$ such that for all $x, y \in \mathbb{R}$ with $|x|^p + |y|^p = 2$:
\[
(x-y)^2 \leq K_p\big(4 - (x+y)^2\big).
\]

\subsection*{Solution}
For $0 < \delta < 1$ we show there is $K_{p,\delta} > 0$ with
\[
f(x,y) = \frac{(x-y)^2}{4 - (x+y)^2} \leq K_{p,\delta}
\]
on $D_\delta = \{(x,y) : |x-y| \geq \delta, |x|^p + |y|^p = 2\}$.

$D_\delta$ is compact; need $4 - (x+y)^2 \neq 0$. Otherwise $|x+y|=2$ and $|(x+y)/2|^p = 1$. Since $p > 1$, $t \mapsto |t|^p$ is strictly convex, so $|(x+y)/2|^p < (|x|^p + |y|^p)/2 = 1$ when $x \neq y$ — contradiction.

If $x, y$ have opposite signs then $|x-y| \geq \max(|x|,|y|) \geq 1 > \delta$, so $(x,y) \in D_\delta$. Thus WLOG $x, y > 0$ with $x^p + y^p = 2$. Set $x = 1 + t$. Then
\[
y = (2 - (1+t)^p)^{1/p} = 1 - t - (p-1)t^2 + o(t^2).
\]
Hence $(x-y)^2 = 4t^2 + o(t^2)$ and $4 - (x+y)^2 = 4(p-1)t^2 + o(t^2)$. So for some $\delta_p > 0$, if $|t| < \delta_p$:
\[
(x-y)^2 < 5t^2 = \frac{5}{3(p-1)} \cdot 3(p-1)t^2 < \frac{5}{3(p-1)}\big(4 - (x+y)^2\big). \quad (*)
\]
This holds when $|x - 1| < \delta_p$; symmetrically when $|y - 1| < \delta_p$.

Finally if $|x-1| \geq \delta_p$ and $|y-1| \geq \delta_p$: since $\max(x,y) \geq 1$ (as $((x+y)/2)^p \leq 1$ gives $x+y \leq 2$), say $x - 1 \geq \delta_p$, then $x - y \geq 2(x-1) \geq 2\delta_p$, giving the bound.

\begin{center}\textbf{Second day}\end{center}

\section*{Problem 1 (Day 2)}
Let $A$ be a $3\times 3$ real matrix such that $Au \perp u$ for every $u \in \mathbb{R}^3$. Prove:
(a) $A^\top = -A$;
(b) there exists $v \in \mathbb{R}^3$ with $Au = v \times u$ for every $u$.

\subsection*{Solution}
(a) Set $A = (a_{ij})$. Using $(Au, u) = 0$ with $u = e_k$ (standard basis), we get $a_{kk} = 0$. With $u = e_k + e_m$:
\[
a_{kk} + a_{km} + a_{mk} + a_{mm} = 0 \implies a_{km} = -a_{mk}.
\]

(b) Set $v_1 = -a_{23}$, $v_2 = a_{13}$, $v_3 = -a_{12}$. Then
\[
Au = (v_2 u_3 - v_3 u_2,\ v_3 u_1 - v_1 u_3,\ v_1 u_2 - v_2 u_1)^\top = v \times u.
\]

\section*{Problem 2 (Day 2)}
Let $b_0 = 1$, $b_n = 2 + \sqrt{b_{n-1}} - 2\sqrt{1 + \sqrt{b_{n-1}}}$. Calculate $\sum_{n=1}^\infty b_n 2^n$.

\subsection*{Solution}
Put $a_n = 1 + \sqrt{b_n}$ for $n \geq 0$. Then $a_0 = 2$ and
\[
a_n = 1 + \sqrt{1 + a_{n-1} - 2\sqrt{a_{n-1}}} = 1 + |\sqrt{a_{n-1}} - 1| = \sqrt{a_{n-1}},
\]
so $a_n = 2^{2^{-n}}$. Then $b_n = (a_n - 1)^2$. Compute
\[
\sum_{n=1}^N b_n 2^n = \sum_{n=1}^N (a_n - 1)^2 2^n = \sum_{n=1}^N [a_n^2 \cdot 2^n - a_n 2^{n+1} + 2^n].
\]
Using $a_n^2 = a_{n-1}$, this telescopes:
\[
\sum_{n=1}^N [(a_{n-1} - 1) 2^n - (a_n - 1) 2^{n+1}] = (a_0 - 1) \cdot 2 - (a_N - 1) 2^{N+1} = 2 - 2 \cdot \frac{2^{2^{-N}} - 1}{2^{-N}}.
\]
As $N \to \infty$, with $x = 2^{-N} \to 0$:
\[
\sum_{n=1}^\infty b_n 2^n = 2 - 2\lim_{x\to 0}\frac{2^x - 1}{x} = 2 - 2\ln 2.
\]

\section*{Problem 3 (Day 2)}
Let all roots of $P(z)$ (degree $n$) lie on the unit circle. Prove all roots of $2zP'(z) - nP(z)$ do too.

\subsection*{Solution}
WLOG $P$ is monic: $P(z) = \prod_{k=1}^n (z - \alpha_k)$, $|\alpha_k| = 1$. Let $\widetilde P(z) = 2zP'(z) - nP(z)$. Expanding,
\[
\widetilde P(z) = \sum_{k=1}^n \prod_{j\neq k}(z - \alpha_j) \cdot (z + \alpha_k) - \text{(via identity)}.
\]
Equivalently,
\[
\frac{\widetilde P(z)}{P(z)} = \sum_{k=1}^n \frac{z + \alpha_k}{z - \alpha_k}.
\]
For $|\alpha| = 1$, $z \neq \alpha$:
\[
\mathrm{Re}\,\frac{z+\alpha}{z-\alpha} = \frac{|z|^2 - |\alpha|^2}{|z-\alpha|^2} = \frac{|z|^2 - 1}{|z-\alpha|^2}.
\]
Hence $\mathrm{Re}(\widetilde P/P) = \sum \frac{|z|^2 - 1}{|z - \alpha_k|^2}$. If $|z| \neq 1$, this is nonzero, so $\widetilde P(z) \neq 0$.

\section*{Problem 4 (Day 2)}
(a) Prove for every $\varepsilon > 0$ there exist $n$ and reals $\lambda_1, \ldots, \lambda_n$ with
\[
\max_{x\in[-1,1]}\left|x - \sum_{k=1}^n \lambda_k x^{2k+1}\right| < \varepsilon.
\]
(b) Prove for every odd continuous $f$ on $[-1,1]$ and every $\varepsilon > 0$, there exist $n$ and $\mu_1, \ldots, \mu_n$ with
\[
\max_{x\in[-1,1]}\left|f(x) - \sum_{k=1}^n \mu_k x^{2k+1}\right| < \varepsilon.
\]

\subsection*{Solution}
(a) Pick $n$ with $(1 - \varepsilon^2)^n \leq \varepsilon$. Then $|x(1 - x^2)^n| < \varepsilon$ on $[-1,1]$. Taking $\lambda_k = (-1)^{k+1}\binom{n}{k}$,
\[
x - \sum_{k=1}^n \lambda_k x^{2k+1} = x\sum_{k=0}^n (-1)^k\binom{n}{k} x^{2k} = x(1-x^2)^n.
\]

(b) By Weierstrass, there is $p \in \Pi_m$ with $\max |f - p| < \varepsilon/2$. Let $q(x) = \tfrac{1}{2}(p(x) - p(-x))$, the odd part. Then
\[
|f(x) - q(x)| \leq \tfrac{1}{2}|f(x) - p(x)| + \tfrac{1}{2}|f(-x) - p(-x)| < \varepsilon/2. \quad (1)
\]
Write $q(x) = b_0 x + \sum_{k=1}^m b_k x^{2k+1}$. If $b_0 = 0$, done by (1). Else apply (a) with $\varepsilon/(2|b_0|)$:
\[
\max\Big|b_0 x - \sum_{k=1}^n b_0 \lambda_k x^{2k+1}\Big| < \varepsilon/2. \quad (2)
\]
Combining (1) and (2) gives the result with $\max(n,m)$.

\section*{Problem 5 (Day 2)}
(a) Prove every function of the form
\[
f(x) = \frac{a_0}{2} + \cos x + \sum_{n=2}^N a_n \cos(nx)
\]
with $|a_0| < 1$ has both positive and negative values on $[0, 2\pi)$.

(b) Prove $F(x) = \sum_{n=1}^{100} \cos(n^{3/2} x)$ has at least $40$ zeros on $(0, 1000)$.

\subsection*{Solution}
(a) Consider $\int_0^{2\pi} f(x)(1 \pm \cos x)\,dx = \pi(a_0 \pm 1)$. If $f \geq 0$, then $a_0 \geq 1$; if $f \leq 0$, $a_0 \leq -1$. Both contradict $|a_0|<1$.

(b) We show: for each integer $N$, each real $h \geq 24$ and each real $y$, $F_N(x) = \sum_{n=1}^N \cos(xn^{3/2})$ changes sign on $(y, y+h)$. The claim follows.

Consider
\[
I_1 = \int_y^{y+h} F_N(x)\,dx, \quad I_2 = \int_y^{y+h} F_N(x)\cos x\,dx.
\]
If $F_N$ has constant sign on $(y, y+h)$, then $|I_2| \leq |I_1|$. We show $|I_2| > |I_1|$.

Since $|\int_y^{y+h}\cos(\alpha x)dx| \leq 2/|\alpha|$,
\[
|I_1| \leq 2\sum_{n=1}^N n^{-3/2} < 2\Big(1 + \int_1^\infty t^{-3/2}dt\Big) = 6. \quad (1)
\]
For $I_2$:
\[
I_2 = \int_y^{y+h}\cos x \cos(xn^{3/2})dx \quad \text{summed} = \tfrac{1}{2}\int_y^{y+h}(1 + \cos 2x)dx + \tfrac{1}{2}\sum_{n\geq 2}\int\big(\cos(x(n^{3/2}-1)) + \cos(x(n^{3/2}+1))\big)dx = \tfrac{h}{2} + \Delta,
\]
where
\[
|\Delta| \leq \tfrac{1}{2} + 2\sum_{n=2}^N \frac{1}{n^{3/2} - 1}.
\]
Using $n^{3/2} - 1 \geq \tfrac{2}{3}n^{3/2}$ for $n \geq 3$:
\[
|\Delta| \leq \tfrac{1}{2} + \tfrac{2}{2\sqrt 2 - 1} + 3\int_2^\infty t^{-3/2}dt < 6.
\]
Hence $|I_2| > h/2 - 6$. With $h \geq 24$: $|I_2| > 6 \geq |I_1|$.

\section*{Problem 6 (Day 2)}
Suppose $\{f_n\}$ are continuous on $[0,1]$ with
\[
\int_0^1 f_m f_n\,dx = \delta_{mn}, \qquad \sup\{|f_n(x)| : x \in [0,1], n \in \mathbb{N}\} < +\infty.
\]
Show no subsequence $f_{n_k}$ converges pointwise on $[0,1]$.

\subsection*{Solution}
By adding functions $\{g_m\}$ if necessary, WLOG $\{f_n\}$ generates $L^2[0,1]$. Suppose $f_{n_k}(x) \to f(x)$ for all $x$. By boundedness and dominated convergence, for fixed $m$:
\[
0 = \int_0^1 f_m f_{n_k} \to \int_0^1 f_m f.
\]
So $\int f_m f = 0$ for all $m$, hence $f = 0$ a.e. But also $1 = \int f_{n_k}^2 \to \int f^2 = 0$, contradiction.

\end{document}
